Deep-learning inference can run on multicore CPUs, on GPUs, or on both, but GPU speedups are often quoted for the kernels alone. We compare serial C++, OpenMP, CUDA and a hybrid OpenMP--CUDA software pipeline on the same MNIST network (784$\to$256$\to$10, 10{,}000 images, batches of 1 to 4096). Each configuration is measured \val{Sp:runs} times on a node with a Tesla K40m GPU (CUDA 11.7) and 24 CPU threads, and the CPU codes are also run on an Apple M4. All implementations agree with the serial reference (largest logit difference \val{Sp:val:cuda:maxerr}, identical predictions). With 24 threads OpenMP is \val{Sp:omp:speedup:4096:24}$\times$ faster than serial at batch 4096 (efficiency \val{Sp:omp:eff:4096:24}). CUDA is \val{Sp:cuda:speedup:4096}$\times$ faster end-to-end, although its kernels alone are \val{Sp:cuda:kspeedup:4096}$\times$ faster, and for single images it is slower than serial. Copies take at most \val{Sp:cuda:tshare:1}\,\% of the CUDA time, whereas single-threaded preprocessing grows to \val{Sp:cuda:preshare:4096}\,\%. The pipelined hybrid needs \val{Sp:hyb:best:e2e:4096}\,ms per pass at batch 4096, \val{Sp:hyb:vscuda:4096}$\times$ faster than CUDA-only and \val{Sp:hyb:vscpu:4096}$\times$ faster than the best OpenMP configuration, but OpenMP is faster for batches of 1 and 8.
